% GENERATED FILE. Edit canonical lessons, then run npm run generate:books.
% canonical-source-hash: fnv1a64:22e0183a37f8b99c
% canonical-lessons: ZH-C506-yuedui, ZH-C506-yanzou, ZH-C506-xiaotiqin, ZH-C506-jingju, ZH-C506-gemi

\chapter{Music on Stage}
\label{ch:zh-music-on-stage}
\hlchaptermodality{506}

\begin{chapteropening}
\textbf{By the end of this chapter:} \emph{I can talk about a band, performing on an instrument, the violin, Beijing opera and fans.}

Complete the last lesson of chapter 506: I can talk about a band, performing on an instrument, the violin, Beijing opera and fans.
\end{chapteropening}

\section[yuèduì]{yuèduì — band (of musicians)}

\label{lesson:ZH-C506-yuedui}



\begin{quote}
Before the new one: say the Mandarin for a stage (in a theatre), then the Mandarin for to rehearse.
\end{quote}

\subsection*{You'll want to know: yuèduì}
\emph{yuèduì} — \textquotedblleft{}band (of musicians)\textquotedblright{}.

\begin{quote}
\emph{wǒ zài yuèduì tán le sān nián jítā} — I played guitar in a band for three years.
\end{quote}

Its characters come later; for now it is a word for the ear.

\begin{cousinweb}
\emph{yuèduì} is built from \zh{乐} 'music' and \emph{duì} 'a troop', both inherited from Middle Chinese, \emph{ngaewk} and \emph{dwojH} (Guangyun, in Baxter's transcription), which glosses the first 'music' and the second 'a band of people'.
\end{cousinweb}

\begin{grammarlens}[title={What you've built}]
One more word for saying how it was.
\end{grammarlens}

\subsection*{Your turn}
\begin{itemize}
\raggedright
  \item \emph{Say it:} \emph{yuèduì}
  \item \emph{Say it:} \emph{yuèduì}, once more
  \item \emph{Say it:} \emph{wǒ zài yuèduì tán le sān nián jítā} — I played guitar in a band for three years.
  \item \emph{Recall:} say the Mandarin for a stage (in a theatre), then the Mandarin for to rehearse, then say \emph{yuèduì} again
\end{itemize}

\begin{tcolorbox}[breakable,colback=teal!4,colframe=teal!35!black,title={Before you move on}]
What does yuèduì mean? (\textquotedblleft{}Band (of musicians)\textquotedblright{}.) Say it once more.
\end{tcolorbox}
\section[yǎnzòu]{yǎnzòu — to perform on an instrument}

\label{lesson:ZH-C506-yanzou}



\begin{quote}
Before the new one: say the Mandarin for to rehearse, then the Mandarin for a band (of musicians).
\end{quote}

\subsection*{You'll want to know: yǎnzòu}
\emph{yǎnzòu} — \textquotedblleft{}to perform on an instrument\textquotedblright{}.

\begin{quote}
\emph{tā yǎnzòu de hěn hǎo} — He plays beautifully.
\end{quote}

Its characters come later; for now it is a word for the ear.

\begin{cousinweb}
\emph{yǎnzòu} is built from \emph{yǎn} 'to perform' and \emph{zòu} 'to present', both inherited from Middle Chinese, \emph{yenX} and \emph{tsuwH} (Guangyun, in Baxter's transcription), which glosses the first 'to spread wide' and the second 'to offer up'; the sense 'to play music' is not given there.
\end{cousinweb}

\begin{grammarlens}[title={What you've built}]
One more word for saying how it was.
\end{grammarlens}

\subsection*{Your turn}
\begin{itemize}
\raggedright
  \item \emph{Say it:} \emph{yǎnzòu}
  \item \emph{Say it:} \emph{yǎnzòu}, once more
  \item \emph{Say it:} \emph{tā yǎnzòu de hěn hǎo} — He plays beautifully.
  \item \emph{Recall:} say the Mandarin for to rehearse, then the Mandarin for a band (of musicians), then say \emph{yǎnzòu} again
\end{itemize}

\begin{tcolorbox}[breakable,colback=teal!4,colframe=teal!35!black,title={Before you move on}]
What does yǎnzòu mean? (\textquotedblleft{}To perform on an instrument\textquotedblright{}.) Say it once more.
\end{tcolorbox}
\section[xiǎotíqín]{xiǎotíqín — violin}

\label{lesson:ZH-C506-xiaotiqin}



\begin{quote}
Before the new one: say the Mandarin for a band (of musicians), then the Mandarin for to perform on an instrument.
\end{quote}

\subsection*{You'll want to know: xiǎotíqín}
\emph{xiǎotíqín} — \textquotedblleft{}violin\textquotedblright{}.

\begin{quote}
\emph{wǒ xué le bā nián xiǎotíqín} — I've been learning the violin for eight years.
\end{quote}

Its characters come later; for now it is a word for the ear.

\begin{cousinweb}
\emph{xiǎotíqín} is built from \zh{小} 'small', \emph{tí} 'to lift' and \emph{qín} 'a zither', all inherited from Middle Chinese, \emph{sjewX}, \emph{dej} and \emph{gim} (Guangyun, in Baxter's transcription), which glosses the second 'to carry by hand' and the third as a stringed instrument; the name for the violin is modern.
\end{cousinweb}

\begin{grammarlens}[title={What you've built}]
One more word for saying how it was.
\end{grammarlens}

\subsection*{Your turn}
\begin{itemize}
\raggedright
  \item \emph{Say it:} \emph{xiǎotíqín}
  \item \emph{Say it:} \emph{xiǎotíqín}, once more
  \item \emph{Say it:} \emph{wǒ xué le bā nián xiǎotíqín} — I've been learning the violin for eight years.
  \item \emph{Recall:} say the Mandarin for a band (of musicians), then the Mandarin for to perform on an instrument, then say \emph{xiǎotíqín} again
\end{itemize}

\begin{tcolorbox}[breakable,colback=teal!4,colframe=teal!35!black,title={Before you move on}]
What does xiǎotíqín mean? (\textquotedblleft{}Violin\textquotedblright{}.) Say it once more.
\end{tcolorbox}
\section[Jīngjù]{Jīngjù — Beijing opera}

\label{lesson:ZH-C506-jingju}



\begin{quote}
Before the new one: say the Mandarin for to perform on an instrument, then the Mandarin for a violin.
\end{quote}

\subsection*{You'll want to know: Jīngjù}
\emph{Jīngjù} — \textquotedblleft{}Beijing opera\textquotedblright{}.

\begin{quote}
\emph{nǐ kàn guo Jīngjù ma} — Have you ever seen Beijing opera?
\end{quote}

Its characters come later; for now it is a word for the ear.

\begin{cousinweb}
\emph{Jīngjù} is built from \emph{Jīng}, short for the capital Beijing, and \emph{jù} 'a play', both inherited from Middle Chinese, \emph{kjaeng} and \emph{gjaek} (Guangyun, in Baxter's transcription), which glosses the first 'great' and the second 'to increase' and 'hard'; the sense 'a play' is not given there.
\end{cousinweb}

\begin{grammarlens}[title={What you've built}]
One more word for saying how it was.
\end{grammarlens}

\subsection*{Your turn}
\begin{itemize}
\raggedright
  \item \emph{Say it:} \emph{Jīngjù}
  \item \emph{Say it:} \emph{Jīngjù}, once more
  \item \emph{Say it:} \emph{nǐ kàn guo Jīngjù ma} — Have you ever seen Beijing opera?
  \item \emph{Recall:} say the Mandarin for to perform on an instrument, then the Mandarin for a violin, then say \emph{Jīngjù} again
\end{itemize}

\begin{tcolorbox}[breakable,colback=teal!4,colframe=teal!35!black,title={Before you move on}]
What does Jīngjù mean? (\textquotedblleft{}Beijing opera\textquotedblright{}.) Say it once more.
\end{tcolorbox}
\section[gēmí]{gēmí — fan (of a singer)}

\label{lesson:ZH-C506-gemi}



\begin{quote}
Before the new one: say the Mandarin for a violin, then the Mandarin for Beijing opera.
\end{quote}

\subsection*{You'll want to know: gēmí}
\emph{gēmí} — \textquotedblleft{}fan (of a singer)\textquotedblright{}.

\begin{quote}
\emph{wǒ shì nàge gēshǒu zuì dà de gēmí} — I'm that singer's biggest fan.
\end{quote}

Its characters come later; for now it is a word for the ear.

\begin{cousinweb}
\emph{gēmí} is built from \zh{歌} 'a song' and \emph{mí} 'to be led astray', both inherited from Middle Chinese, \emph{ka} and \emph{mej} (Guangyun, in Baxter's transcription), which glosses the second 'to be confused'.
\end{cousinweb}

\begin{grammarlens}[title={What you've built}]
One more word for saying how it was.
\end{grammarlens}

\subsection*{Your turn}
\begin{itemize}
\raggedright
  \item \emph{Say it:} \emph{gēmí}
  \item \emph{Say it:} \emph{gēmí}, once more
  \item \emph{Say it:} \emph{wǒ shì nàge gēshǒu zuì dà de gēmí} — I'm that singer's biggest fan.
  \item \emph{Recall:} say the Mandarin for a violin, then the Mandarin for Beijing opera, then say \emph{gēmí} again
\end{itemize}

\begin{tcolorbox}[breakable,colback=teal!4,colframe=teal!35!black,title={Before you move on}]
What does gēmí mean? (\textquotedblleft{}Fan (of a singer)\textquotedblright{}.) Say it once more.
\end{tcolorbox}
